\documentclass[12pt]{article}
\usepackage[utf8]{inputenc}

\usepackage{mystyle}
\usepackage{parskip}
\usepackage{float}
\usepackage[normalem]{ulem}

% Comic Sans
% \usepackage{sansmathfonts}
% \usepackage{fontspec}
% \setmainfont{Comic Sans MS}

% \usepackage{fontspec}
% \setmainfont{Wingdings-Regular.ttf}

% \usepackage{libertine}

% fancy header
\usepackage{fancyhdr}
\pagestyle{fancy}
\rhead{Joseph Sullivan}

\title{Koszul Duality}
\author{Joseph Sullivan}
\date{January 2025}

\newcommand{\reg}{\mathrm{reg}}
\newcommand{\sing}{\mathrm{sing}}
\newcommand{\Pic}{\mathrm{Pic}}
\newcommand{\sgn}{\mathrm{sgn}}
\newcommand{\FS}{\mathrm{FS}}
\newcommand{\Alb}{\mathrm{Alb}}
\newcommand{\CMan}{\text{$\C$-$\mathsf{Man}$}}

\newcommand{\ev}{\mathrm{ev}}

\newcommand{\Hol}{\mathrm{Hol}}

\newcommand{\Sym}{\mathrm{Sym}}
\newcommand{\Tor}{\mathrm{Tor}}

\DeclareMathOperator{\Weil}{Weil}
\DeclareMathOperator{\Div}{Div}
\DeclareMathOperator{\CDiv}{CDiv}
\DeclareMathOperator{\codim}{codim}
\DeclareMathOperator{\height}{height}
\DeclareMathOperator{\Ass}{Ass}
\renewcommand{\div}{\operatorname{div}}

\DeclareMathOperator{\ord}{ord}
\DeclareMathOperator{\Cl}{Cl}

\begin{document}

\maketitle

\tableofcontents

\section{An Example}

\subsection{Setup}
We start with a computation. Let $k$ be a field, $V$ a $k$-vector space, and $S=\Sym^\bullet V$ the symmetric algebra on $V$ (which is just the $k$-polynomial ring on a basis $E^1,\dots,E^n$ of $V$).

The ring $S$ is graded, with the module $k=S_0\cong S/S_+$. Our task is to compute the Ext algebra
\[\Ext_S^\bullet(k,k),\]
by which we mean the direct sum of all $\Ext_S^i(k,k)$ equipped with multiplication by the Yoneda product (which is either thought of as composition in the derived category of $\Ext_S^i(k,k)=\Hom_{D^b(S)}(k,k[i])$ or as concatenation of extensions).

To compute in the derived category, we will want to replace $k$ (thought of as a complex concentrated in degree 0) with a quasi-isomorphic object, a projective resolution given by the Koszul complex:
\[0\to \bigwedge^n V \otimes_k S \to \cdots \to \bigwedge^1 V \otimes_k S \to \bigwedge^0 V \otimes_k S \to 0,\]
which we will denote by $P^\bullet$ (where $P^i=\bigwedge^{-i} V \otimes_k S$).

When we say that $P^\bullet$ is a projective resolution of $k$, we mean that each term is projective, and we have a map $P^\bullet \xrightarrow{\eps} k$ sending $\bigwedge^0 V \otimes_k S \to k$ (by evaluating $s\in S$ at $0$) which is a quasi-isomorphism. This will turn out to mean that $P^\bullet$ is an exact complex, except at $P^0$, where the map $P^{-1}\to P^0$ has cokernel $k$.

Now, with this in mind,
\[\Hom_{D^b(S)}(k,k[i]) \cong \Hom_{D^b(S)}(P^\bullet,k[i]) \cong \Hom_{D^b(S)}(P^\bullet, P^\bullet[i]),\]
where we get the isomorphism by pre/post composing a morphism $k\to k[i]$ with the isomorphism (in $D^b(S)$) $P^\bullet \cong k$.

Now, it is a theorem on derived categories that
\[\Hom_{D^b(S)}(k, k) \cong \Hom_{K^b(S)}(P^\bullet, k[i]) \cong \Hom_{K^b(S)}(P^\bullet, P^\bullet[i]),\]
that is, when the domain object is a complex consisting of projective objects, morphisms in the derived category are those coming from the homotopy category.

Now, to compute the $\Ext$ algebra, we will start with the more straightforward computation
\[\bigwedge^i V^\vee \cong \Hom_{K^b(S)}(P^\bullet,k[i]),\]
but then we will work out explicitly how this is isomorphic to $\Hom_{K^b(S)}(P^\bullet,P^\bullet[i])$, because it will allow us to compute multiplication. Finally, by using these isomorphisms, we will see that in the diagram
\begin{cd}
    \Hom_{K^b(S)}(P^\bullet,P^\bullet[i]) \otimes \Hom_{K^b(S)}(P^\bullet[i],P^\bullet[i+j]) \dar \rar & \Hom_{K^b(S)}(P^\bullet,P^\bullet[i+j]) \dar\\
    \bigwedge^i V^\vee \otimes \bigwedge^j V^\vee \rar & \bigwedge^{i+j} V^\vee
\end{cd}
the bottom arrow is given by the familiar multiplication
\[(f_1\wedge \cdots \wedge f_i) \wedge (g_1 \wedge \cdots g_j) = f_1 \wedge \cdots \wedge f_i \wedge g_1 \wedge \cdots \wedge g_j.\]

\subsection{Computation}

\subsubsection*{Step 1.}
We write down an isomorphism
\[\bigwedge^i V^\vee \cong \Hom_{K^b(S)}(P^\bullet, k[i]).\]
To get the isomorphism, we recognize that the data is given by a chain map $P^\bullet \to k[i]$ is a $S$-module morphism
\[\bigwedge^i V \otimes_k S \longrightarrow k\]
that precomposes to 0 under the boundary map $\bigwedge^{i+1} V \otimes_k \to \bigwedge^i V \otimes_k S$, but this extra condition is vacuous because any $S$-module morphism to $k=S/S_+$ will kill
\[S_+ (\bigwedge^i V \otimes_k S)\]
which contains the image of the boundary map.

So, the data is exactly
\[\Hom_S(\bigwedge^i V \otimes_k S, k),\]
and here we can write down a morphism
\begin{align*}
    \bigwedge^i V^\vee &\longrightarrow \Hom_S(\bigwedge^i V \otimes_k S, k)\\
    f_1\wedge\cdots\wedge f_i &\longmapsto (v_1\wedge \cdots \wedge v_i \mapsto \det [f_p(v_q))]_{pq}
\end{align*}
This is an isomorphism, since one can check this is a surjection by picking a basis $e_1,\dots,e_n$ of $V$ (and $e_1^\vee,\dots,e_n^\vee$ of $V^\vee$ the dual basis) and recognizing that for $I\subseteq \{1,\dots,n\}$ with $|I|=i$,
\[e_I^\vee \mapsto e^I \otimes (S\to k)\]
which span the target, and then both vector spaces have the same dimension.

\subsubsection*{Step 2.}
Now, we want to carry over this isomorphism
\[\bigwedge^i V^\vee \cong \Hom_{K^b(S)}(P^\bullet, k[i]) \cong \Hom_{K^b(S)}(P^\bullet,P^\bullet[i]).\]
To do so, given $\omega \in \bigwedge^i V^\vee$, we want to produce a morphism $P^\bullet\to P^\bullet[i]$, such that after composing with $P^\bullet[i] \xrightarrow{\eps[i]} k[i]$, we get the morphism described in Step 1. The general theory tells us that
\[\Hom_{K^b(S)}(P^\bullet,P^\bullet[i]) \xrightarrow{\eps[i]_*} \Hom_{K^b(S)}(P^\bullet, k[i])\]
is an isomorphism, so as soon as we find any chain map realizing $\omega$, we are done (in the sense that we have realized the homotopy class of the chain map).

So, we need to build a chain map as in the diagram.
\begin{equation}\label{extend-to-koszul}
\begin{tikzcd}
    \cdots \rar & \bigwedge^{i+1} V \otimes_k S \rar \dar[dashed] & \bigwedge^i V \otimes_k S \rar \dar[dashed] & \bigwedge^{i-1} V \otimes_k S \rar \dar[dashed] & \cdots\\
    \cdots \rar & \bigwedge^1 V \otimes_k S \rar & \bigwedge^0 V \otimes_k S \rar \dar["\eps"] & 0 \rar & \cdots\\
    & & k & &
\end{tikzcd}
\end{equation}


We can understand the maps in the Koszul complex by taking the embeddings
\begin{align*}
    \bigwedge^{r+1} V &\longrightarrow \bigwedge^{r} V \otimes_k V\\
    v_1\wedge \cdots \wedge v_{r+1} &\longmapsto \sum_{\ell=1}^{r+1} (-1)^\ell (v_1 \wedge \cdots \wedge \hat{v_\ell} \wedge \cdots \wedge v_{r+1}) \otimes v_\ell.
\end{align*}
tensoring them with $S$, and composing with the multiplication (tensored with $\bigwedge^r V$)
\[V \otimes_k S \longrightarrow S.\]
In this way, our goal is to get commuting squares (when $j\geq i$)
\begin{equation}\label{koszul-boundary}
\begin{tikzcd}
    \bigwedge^{j+1} V \rar \dar[dashed, "\iota_\omega^{j+1}"] & \bigwedge^{j} V \otimes V \dar[dashed, "\iota_\omega^j \otimes \id_V"]\\
    \bigwedge^{j-i+1} V \rar &  \bigwedge^{j-i} V \otimes V
\end{tikzcd}
\end{equation}
which we will then tensor with $S$ and compose with the $V$ multiplication on $S$ to get our desired chain map. In this way, we have reduced the problem to linear algebra, where our task is to find a certain natural map $\bigwedge^{j} V \to \bigwedge^{j-i} V$ induced by $\omega \in \bigwedge^i V^\vee$.

This will be some sort of interior multiplication, a map
\[\bigwedge^i V^\vee \otimes \bigwedge^j V \longrightarrow \bigwedge^{j-i} V.\]
Writing down this morphism takes some combinatorics. Denote $J=\{1,\dots,j\}$, and for any subset $I=\{i_1<\cdots<i_i\}\subseteq J$, we can define
\[\sgn_J(I) = \sum_{\ell=1}^i (i_\ell-\ell) \mod{2} \mapsto \{-1,1\}\]
which measures how many moves it takes (while keep $I\setminus J$ in order) to get from $J$ to $\{1,\dots,j\} \subseteq I$. Now, we define
\begin{align*}
    \bigwedge^i V^\vee \otimes \bigwedge^j V &\longrightarrow \bigwedge^{j-i} V\\
    (f_1\wedge \cdots \wedge f_i) \otimes (v_1\wedge \cdots \wedge v_j) &\longmapsto \sum_{\substack{I\subseteq J\\ |I|=i}} (\sgn_J I) \det[f_p(v_{i_q})]_{pq} v_{J\setminus I}.
\end{align*}
We need to check that this is well-defined, in the sense that we have really defined a set function
\[(V^\vee)^{\times i} \times V^{\times j} \to \bigwedge^{j-i} V,\]
and we want to see that it factors through the projection to the tensor of wedges. We will need to check multilinearity and alternating properties.

First, it is straightforward to check that this is multilinear in $f_1,\dots,f_i,v_1,\dots,v_j$. Next, to check that it is alternating on the first $i$ factors or last $j$ factors, we can check without loss of generality that when either $f_1=f_2$ or $v_1=v_2$, we get 0.

When $f_1=f_2$, every term of the sum is zero, since we are evaluating determinants of matrices with repeat rows. When $v_1=v_2$, the picture is more complicated. We do some case work:

\textbf{Summands with $\{1,2\}\subseteq I$.} In this case, the determinant will have repeat columns, so each of these summands are 0.

\textbf{Summands with $\{1,2\}\cap I = \emptyset$.} In this case, $v_{J\setminus I}$ will have a wedge factor $v_1\wedge v_2$, which is 0, so each of these summands are 0.

\textbf{Summands with $|\{1,2\}\cap I|=1$.} Here, the summands need not be 0, but we want to argue that they ``cancel each other out''. We have a bijection
\[\{I\subseteq J \mid 1\in I, 2\not\in I\} \cong \{I \subseteq J \mid 2\in I, 1\not\in J\},\]
just by sending the flipping first index $i_1$ from $1$ to $2$ or vice versa. Denote this by $I\mapsto -I$, which satisfies the $\sgn_J -I = -\sgn_J I$. This will mean that the summand corresponding to $I$ and that corresponding to $-I$ will have opposite signs and annihilate each other.

Thus, we have succeeded in defining our interior multiplication map. Specializing to $i=j$, we get
\begin{align*}
    \bigwedge^i V^\vee \otimes \bigwedge^i V &\longrightarrow \bigwedge^0 V \cong k\\
    (f_1\wedge \cdots \wedge f_i) \otimes (v_1 \wedge \cdots \wedge v_j) &\longmapsto \det[f_p(v_q)],
\end{align*}
so, assuming this interior multiplication gives us an actual chain map, this specialization tells us that it lifts the chain map $P^\bullet \to k[i]$ induced by an $\omega\in \bigwedge^i V^\vee$.

Now, we show that the interior multiplication actually gives a chain map, i.e., that it actually commutes with Koszul boundaries, so we will check that a beefed up version of diagram (\ref{koszul-boundary}) commutes:
\begin{cd}[column sep=2cm]
    \bigwedge^i V^\vee \otimes \bigwedge^{j+1} V \rar["\id_{\bigwedge^i V^\vee} \otimes d"] \dar["\iota^{j+1}"] & \bigwedge^i V^\vee \otimes \bigwedge^j V \otimes V \dar["\iota^j \otimes \id_V"] \\
    \bigwedge^{j-i+1} V \rar & \bigwedge^{j-i} V \otimes V
\end{cd}
We check commutativity by evaluating these maps on basis elements. Let $e^\vee_{I'} \otimes e_{\tilde{J}} \in \bigwedge^i V^\vee \otimes \bigwedge^{j+1} V$. Going right, we get
\[e^\vee_{I'} \otimes \left(\sum_{\ell=1}^{j+1}(-1)^\ell e_{\tilde{J}\setminus \{\tilde{j}_\ell\}} \otimes e_{\tilde{j}_\ell}\right),\]
which then maps down to
\[\sum_{\ell=1}^{j+1} \sum_{I\subseteq J_\ell=\tilde{J}\setminus \{j_\ell\}} \sgn_{J_\ell}(I) \det[e^\vee_{i'_p}(e_{j_{i_q})}]_{pq} e_{J_\ell \setminus I} \otimes e_{\tilde{j}_\ell}.\]
Now, the determinant factor is given by the Kronecker delta $\delta_{I'I}$, i.e., it is 1 when $I=I'$, and zero otherwise. For each $\ell$, either $\tilde{j}_\ell \in I'$ and then we will not get $I'$ when iterating through $I\subseteq J_\ell$, or $\tilde{j}_\ell \not\in I'$ and we do get $I'$. This will tell us that our term is given by
\[\sum_{\tilde{j}_\ell \not\in I'} \sgn_{J_\ell}(I') e_{J_\ell \setminus I'} \otimes e_{\tilde{j}_\ell}.\]

Now, we compute the other direction. Again starting with $e^\vee_{I'} \otimes e_{\tilde{J}}$, we first go down to get
\[\sum_{I\subseteq \tilde{J}} \sgn_{\tilde{J}}(I) \det[e^\vee_{i'_p}(e_{\tilde{j}_{i_q}})] e_{\tilde{J}\setminus I},\]
but the determinant will only be nonzero when $I=I'$, so we just get the term
\[\sgn_{\tilde{J}}(I') e_{\tilde{J}\setminus I'}.\]
Denote $\hat{J}=\tilde{J}\setminus I'$. This then maps to the right to get
\[\sgn_{\tilde{J}}(I) \sum_{\ell=1}^{j-i+1} (-1)^\ell e_{\hat{J} \setminus \{\hat{j}_\ell\}} \otimes e_{\hat{j}_\ell}.\]
Up to sign this is correct, one should check more carefully to see it is fully correct.

Now, we see that a given $\omega \in \bigwedge^i V^\vee$, we really can construct a lift as in diagram (\ref{extend-to-koszul}), which we denote by $\iota^\bullet_\omega$, by tensoring $S$ with interior multiplication by $\omega$.

\textbf{Step 3.}

Now, we compute the multiplication between $\omega\in \bigwedge^i V^\vee$ and $\eta\in \bigwedge^j V^\vee$, by computing the composition
\[\iota_\eta^{j} \circ \iota_\omega^{i+j}: \bigwedge^{i+j} V \to k,\]
(which maybe should be tensored with $S$) and we check that this corresponds to $\eta \wedge \omega$. We again compute on a basis. Suppose $\omega = e_{I'}^\vee$ and $\eta = e_{J'}^\vee$. If $I',J'$ are not disjoint, we expect multiplication to be 0. So, assume they are not disjoint, and let $K$ be a length $i+j$ subset. Evaluating first $\iota_\omega^{i+j}$ on $e_K$, we only get a nonzero value if $I'\subseteq K$, and in that case we get
\[e_K \mapsto \sgn_K(I') e_{K\setminus I'},\]
but then under $\iota_\eta^{j}$ we get 0 because $J'\not\subseteq K\setminus I'$.

So, now assume $I',J'$ are disjoint. We will similarly expect to get just the indicator function (up to sign) on $e_K$, and indeed this works out---if $K$ doesn't contain both $I',J'$, then one of the computations will be 0. Otherwise, $K=I' \sqcup J'$, and we will get exactly the product of signs.

This tells us $\iota_\eta^{j} \circ \iota_\omega^{i+j} = \iota_{\eta \wedge \omega}^{i+j}$, which tells us that the multiplication in $\Ext_S^\bullet(k,k)$ is exactly that of $\bigwedge^\bullet V^\vee$.

\end{document}